\documentclass[10pt]{article}
\usepackage[utf8]{inputenc}
\usepackage[T1]{fontenc}
\usepackage{amsmath}
\usepackage{amsfonts}
\usepackage{amssymb}
\usepackage[version=4]{mhchem}
\usepackage{stmaryrd}

\begin{document}
была установлена теорема о переменной О. к. п., для к-рой теорема Гильберта является предельным случаем. Иначе говоря, те свойства постоянной О. к. п., к-рые вызывают эффект теоремы Гильберта (нарушение регулярности), и сам әффект возможно понимать в приближенном смысле. Для доказательства было найдено уравнение в случае переменной кривизны, к-рое обобщает уравнение (2). Именно,

\[
k^{2} \frac{\partial}{\partial s_{2}}\left(\frac{1}{K^{2}} \frac{\partial \omega}{\partial s_{1}}\right)=\left\{k^{2}+A+\alpha|\operatorname{grad} \omega|\right\} \sin \omega,
\]

где $\omega$ имеет прежний смысл, $k^{2}=-K$, диффференцирование ведется по дугам асимптотич. линий, $A$ и $\alpha-$ величины, модули к-рых допускают внутреннюю оценку, т. е. оценку в зависимости только от погружаемой метрики. При $K=$ const получается $A=0$, $\alpha=0$ и уравнение (4) превращается в уравнение (2).

В случае переменной кривизны найдено (см. [17]) уравнение

\[
L_{s_{2}} L_{s_{1}}^{*}(\omega)=M \sin \omega,
\]

где $L_{s_{2}}$-линейныї оператор, пропорциональный оператору $^{2}$ диффференцирования по дуге асимптотич. линий второго семейства, $L_{s_{1}}^{*}$ - квазилинейный оператор похожей структуры. Существенно, что функция $M$ при нек-рых условиях на метрику допускает внутреннюю оценку снизу положительным числом, благодаря чему и благодаря строению левой части уравнение (5) сходно с уравнением (2). Поэтому при таких условиях на метрику можно провести рассуждения, сходные с рассуждениями Д. Гильберта (хотя гораздо более сложные), и получить сходный результат. Таким путем впервые было получено обобщение теоремы Гильберта (см. [18]). Что касается условий на метрику, о к-рых только что говорилось, то они заключаются в требовании, чтобы гауссова кривизна $K$ была отгорожена от нуля и медленно изменялась. Если предположить, что $K \leqslant-1$, то медленное изменение кривизны нужно понимать в том смысле, что первые и вторые производные от $K$ по дуге любой геодезической достаточно малы. Достаточные оценки устанавливаются с таким расчетом, чтобы к уравнению (5) были применены указанные выше рассуждения, идущие от Д. Гильберта. Эти оценки (см. [19]) выражены с помощью положительных постоянных, обозначаемых $h$ и $\Delta$. Соответственно, получаемые метрики с медленно изменяющейся отрицательной кривизной наз. $(h, \Delta)$ )м е т р и к а м и (определение см. в [18] или в [19]); при этом метрики постоянной кривизны получаются при $\Delta=+\infty$. Таким образом, полная $(h, \Delta)$-метрика не допускает регулярного изометрич. погружения в трехмерное евклидово пространство.

Вопрос о возможности обобщения теоремы Гильберта в том виде, как он ставился в 30-е гг. (без ограничений характера поведения кривизны, см. [12]), на указанных сейчас путях, т. е. с помощью дифференциальных уравнений теории поверхностей, завершить не удалось. Но решение этого вопроса получилось совсем на другом пути, именно, с помощью тщательного исследования граничных свойств сферич. образа О. к. п. Таким способом доказано (см. [19], [20]), что полная метрика с гауссовой кривизной $K \leqslant$ const $<0$ не допускает регулярного изометрич. погружения в $E_{3}$. Тем самым вопрос решен точно в классич. постановке. Одновременно доказана теорема: в $E_{3}$ на всякой полной регулярной поверхности отрицательной гауссовой кривизны $K$ имеет место равенство

\[
\sup K=0 \text {. }
\]

Здесь достаточна регулярность $C^{2}$. Но в случае $C^{1}, 1$ есть (см. [7]) контрпример (в виде слабо нерегулярной поверхности). Кроме (6) получены также (см. [21]) другие равенства и оценки, универсальные в том смысле, что они относятся ко всем О. к. п. в $E_{3}$. Кроме того, метод доказательства (6) позволил получить теоремы, к-рые дают достаточные признаки, когда отображение плоскости в плоскость с отрицательным якобианом является не только локальным, но и глобальным диф্ффеоморфизмом.

Несмотря на изложенные сейчас результаты, исследование $(h, \Delta)$-метрик и других метрик с медленно изменяющейся кривизной (напр., q-метрик; см. [9]) не потеряло смысла. Дело в том, что в задачах погружения метрик с медленно изменяющейся отрицательной кривизной возникают әффекты, к-рые не имеют места в случае более общих метрик отрицательной кривизны (см., напр., теорему о простой зоне в [9]). Образно можно сказать, что чем медленнее изменяется кривизна, тем «труднее» метрике отрицательной кривизны быть погруженной в $E_{3}$. Для метрики постоянной кривизны это «явление» проявляется особенно заметно (см. усиленную теорему Гильберта в [9]). Что касается диффференциальных уравнений, к-рые отмечались в связй с вопросом о погружении $(h, \Delta)$-метрик, то они суџцественно использовались и в других вопросах. Так, с их помощью установлен [7] ряд сильных теорем о связи между внутренними и внешними свойствами О.к. п.; в частности, теоремы о зависимости внешнеӥ регулярности кривизны О. к. п. от регулярности ее метрики (см. [7]).

Влияние регулярности метрики на внешнюю регулярность О. к. П. представляется особенно интересным, так как О. к. п. входят в нек-рый специальный класс поверхностей, к-рый может быть определен чисто геометрически (интересно, что чисто геометрич. определение допускает и класс выпуклых поверхностей). Поверхности этого класса наз. с е д л о в ы м и. Седловые поверхности в нек-ром смысле - антиподы выпуклых (см. [22] - [25]).

Теорема о невозможности в евклидовом пространстве полной O. к. п. $K \leqslant$ const $<0$ получает весьма усиленную форму в случае поверхности, к-рая однолистно проектируется на плоскость, т. е. может быть представлена как график функции $z=f(x, y)$. В этом случае существуют универсальные оценки протяженности поверхности, и полнота ее становится несущественной. Доказаны теоремы (см. [26]): 1) если на прямоугольник со сторонами $a, b$ регулярно проектируется поверхность с гауссовой кривизной $K \leqslant-\alpha^{2}=$ const $<0$, то неравенство $\sqrt{\alpha} \geqslant C_{1}+\varepsilon, \varepsilon>0$ влечет $a \sqrt{\alpha} \leqslant C_{1}+\frac{C_{2}}{\varepsilon}$, где $C_{1}, C_{2}$ - универсальные постоянные (напр., $C_{1}=4 \pi$, $\left.C_{2}=12 \pi\right) ; 2$ ) если $a, b$ и (для простоты) $\alpha=1$, то $a \leqslant 18,9$.

Пусть дана метрика $d s^{2}$ с отрицательной гауссовой кривизной $K$ и $k=\sqrt{-K}$. Пусть метрика $d s^{2}$ погружена в $E_{3}$ и $l, m, n$ - приведенные коэфффициенты получившейся второй квадратичной формы (приведенные значит удовлетворяюгцие уравнению $l n-m^{2}=-K$ ). В задаче о погружении $l, m, n$ - неизвестные функции. В качестве новых неизвестных функций вводятся т. н. и н в а р и а н ты $\mathbf{P}$ и м ан а:

\[
r=-\frac{m+h}{n}, s=\frac{m-k}{n} .
\]

Тогда основная система уравнений төории поверхностей примет симметричный вид

\[
\left.\begin{array}{l}
r_{x}+s r_{y}=A_{0}+A_{1} r+A_{2} s+A_{3} r^{2}+A_{4} r s+A_{5} r^{2} s \\
s_{x}+r s_{y}=A_{0}+A_{1} s+A_{2} r+A_{3} s^{2}+A_{4} r s+A_{5} s^{2} r
\end{array}\right\}
\]

где $A_{0}, \ldots, A_{5}$ выражаются только через метрику $d s^{2}$.

Выше указывалось, что полная 0 . к. П. с $K \leqslant$ const $<0$ не допускает изометричного регулярного погружения в $E_{0}$. Поэтому полученные результаты о возможности изометрич. погружения двумерных многообразий от-


\end{document}